\section{What was caught}\label{sec:caught}

\Cref{tab:findings} summarises the findings. None was an error in a
machine-checked proof. Three were errors in the pen-and-paper sketch that the
mechanization replaced. Five concerned our first write-up of the
mechanization: three overstatements, and two restrictions that it disclosed
in its notes but left out of the statements of the results. One concerned our
second write-up, and one was a divergence between the implementation and its
specification's algorithm, already reported for another library. Two reviews
contributed: one of the merged mechanization, and one of a draft of this
paper. Separate instances of the AI assistant carried out both
(\cref{sec:intro}).

\begin{table}[!tbp]
\centering
\footnotesize
\begin{tabularx}{\linewidth}{@{}>{\raggedright\arraybackslash}p{0.31\linewidth}>{\raggedright\arraybackslash}p{0.14\linewidth}L@{}}
\toprule
Finding & Found by & Resolution\\
\midrule
\multicolumn{3}{@{}l}{\emph{In the proof sketch (\cref{sec:sketchcaught})}}\\
Composite signature needs a collision-resistant pre-hash & mechanization & \cref{lem:sig} gains $2\,\Adv^{\text{cr}}_{\PH}$, as the draft states~\cite[\S1.3]{Composite}; mutation test\\
X-Wing relying on X25519 needs C2PRI of ML-KEM & mechanization & \cref{lem:kem} gains $\cvp{Pc2pri}$, as in~\cite[Thm.~1]{XWing}; mutation test\\
Theorem~1's reduction guesses an honest party, but its bound has no factor for the guess & second review & the mechanized bound has it: $N_K \cdot \cvp{Psign}$\\[3pt]
\multicolumn{3}{@{}l}{\emph{In our first write-up of the mechanization (\cref{sec:reviewcaught})}}\\
``No session-guessing loss'': bound claimed tighter than the sketch's & first review & retracted: the multi-instance KEM term costs ${\approx}\,n_S^2$ when reduced to single-instance IND-CCA2\\
Non-standard KDF assumptions not flagged as such & first review & named precisely; a random-oracle proof from standard assumptions added\\
``Exactly the X-Wing paper's assumptions'' & first review & qualified: prime-order group here, nominal group there\\
Result statements omitted a disclosed restriction: keys of possible partners withheld from RevealKey & first review & restriction removed: RevealKey on request with partner-based freshness; key uniqueness added\\
Result statements omitted that ephemeral reveal is static & first review & claims say ``static''\\[3pt]
\multicolumn{3}{@{}l}{\emph{In our second write-up (\cref{sec:reviewcaught})}}\\
Mutation coverage overstated: six guard mutants covered seven of twelve guards & second review & one mutant per guard; 22 tests in all\\[3pt]
\multicolumn{3}{@{}l}{\emph{In the implementation (\cref{sec:lowcaught})}}\\
X-Wing raised an error on low-order X25519 points & first review; confirmed by tests & fixed to follow the draft; 63 tests and nine mutants of the fix; reported earlier for another library~\cite{Noble26}; the ML-KEM-1024 profile keeps its check\\
\bottomrule
\end{tabularx}
\caption{Findings, by where the error was and what found it.}
\label{tab:findings}
\end{table}

\subsection{In the proof sketch}\label{sec:sketchcaught}

Before the mechanization, the repository contained a reduction-based proof
sketch in a Bellare--Rogaway-style model~\cite{BR93,CK01}. It had two
theorems for the protocol and two lemmas for the hybrid primitives. Both
theorems were re-proved in a modified model (\cref{sec:freshness}), with
per-role sets of test keys, RevealKey on request, and static rather than
adaptive ephemeral reveal. Both lemmas were wrong as stated, and in both
cases the primary source states the condition that the sketch dropped. The
sketch's Theorem~1 bound also lacked a factor.

\paragraph{The composite signature.} The sketch stated
$\Adv^{\text{euf-cma}}_{\text{composite}} \le \min(\Adv^{\text{euf-cma}}_{\text{ML-DSA}}, \Adv^{\text{euf-cma}}_{\text{EdDSA}})$,
arguing that a composite forgery on a fresh message $M^*$ contains, for each
component, a valid signature on $M^*$ that the component never produced. The
components do not sign $M^*$, though. They sign $M'$, which depends on $M^*$
only through $\PH(M^*)$. If $\PH(M^*) = \PH(M)$ for a message $M$ signed
earlier under the same context, the composite signature on $M$ is also a
valid signature on $M^*$. Modelling the draft's message representative
exactly made the missing assumption unavoidable: the proof needs collision
resistance of $\PH$, and the corrected bound is \cref{lem:sig}. The draft
states this condition itself~\cite[\S1.3]{Composite}; the sketch had dropped
it. The mutant \ident{composite_no_prehash_collision_resistance} confirms that
the proof fails without it. With SHA-512 or SHAKE256 as $\PH$, the extra term
is negligible in practice. Still, the sketch's claim that the composite is as
strong as its stronger component holds only together with a hash
assumption.

\paragraph{X-Wing.} The sketch stated, for both profiles,
$\Adv^{\text{ind-cca}}_{\text{KEM}} \le \min(\Adv^{\text{ind-cca}}_{\text{ML-KEM}}, \Adv^{\text{dh}}_{\text{X25519}}) + \varepsilon_{\text{comb}}$,
on the grounds that the combiner is a dual PRF. For X-Wing relying on
X25519, this misses C2PRI of ML-KEM. The mechanized proof applies that
assumption after the random-oracle step to rule out decapsulation queries
$(\ct_M', \ct_X^*)$ with $\ct_M' \ne \ct_M^*$ that would reach the
challenge's oracle input (\cref{sec:mech:kem}). This matches the X-Wing
paper's classical theorem~\cite[Thm.~1]{XWing}, which bounds the IND-CCA
advantage of its generic construction, of which X-Wing is an instance, by
$2\Delta_N$ plus the strong-DH and C2PRI advantages. The X-Wing paper also
proves that ML-KEM-768 is C2PRI when its hash functions are random
oracles~\cite[Thm.~3]{XWing}, so X-Wing's guarantee survives a break of the
lattice problem. But the ``either leg'' statement is not symmetric: the
X25519 leg needs a property of the ML-KEM leg. The ML-KEM-1024 profile,
which hashes $\ct_M$, needs no security assumption on ML-KEM in its X25519
direction. The sketch's symmetric form does hold for it, once its
unspecified Diffie--Hellman assumption is made precise as strong
Diffie--Hellman with a random oracle. The mutant
\ident{bbfgs_rom_no_ctM_in_combiner} shows that hashing $\ct_M$ is the
reason.

\paragraph{Theorem~1's bound.} The sketch's reduction for Theorem~1 ``plants
the challenge vk at a guessed honest party'' and answers that party's
signing queries with its oracle, but its bound is
$\Adv^{\text{euf-cma}}_{\text{Sig}}$, with no factor for the guess. Guessing
among $n_P$ honest parties loses a factor of $n_P$. The mechanized bound has
the corresponding factor, $N_K \cdot \cvp{Psign}$, through the library's
UF-CMA-with-corruption macro. We had compared the two bounds without
noticing this; the second review caught it.

\subsection{In our write-ups: the reviews}\label{sec:reviewcaught}

After the first version of the mechanization was merged, a separate instance
of the AI assistant reviewed it. It could read the repository and its
history, but not the development conversation. It re-ran every model,
probed the proofs, and found no unsound step and no vacuous query. It did
find places where the write-up claimed more than the models prove. We
re-checked the model-related findings with CryptoVerif before acting on
them.

\paragraph{Tightness.} The first write-up said the mechanized bound had ``no
session-guessing loss'', so that the $n_S^2$ factors of the sketch's bound
were ``an artefact of the sketch's proof strategy''. The first half is true:
CryptoVerif's proof guesses no session. The conclusion is not. The KEM term
$\cvp{Pkemcca}(t, \calls{OI1}, \calls{OA1}+\calls{OR1}, \calls{OI2})$ is a
\emph{multi-key, multi-challenge} IND-CCA2 advantage, with $\calls{OI1}$
keys and up to $\calls{OR1} + \calls{OA1}$ challenges, since one $m_1$ can
be delivered to every responder. Reducing it to single-key, single-challenge
IND-CCA2 generically loses a factor of about keys times
challenges~\cite{BBM00}, that is, about $n_S^2$, the same order as the
sketch. What is better than the sketch is narrower. The signature term is
$N_K \cdot \Adv^{\text{euf-cma}}$, against $n_S^2 \cdot \Adv^{\text{euf-cma}}$
in the sketch's concrete statement (its Theorem~1 omitted the factor
altogether), and the PRF term is linear in the sessions. A tighter overall bound would
need a tight multi-user bound for the KEM, which we do not claim.

\paragraph{Non-standard assumptions not flagged.} The sketch justified ``HKDF
is a PRF keyed by its input keying material'' by HMAC being a
``(dual-)PRF''~\cite{Kra10}. HKDF-Extract keys HMAC with the salt, however,
so keying it through the ikm uses HMAC's message input. That is the swap-PRF
property, which Backendal et al.\ prove only for feasible key
sets~\cite{BBGS23} (\cref{sec:results}). For the ML-KEM-1024 combiner, the
first write-up stated both legs' assumptions without saying that neither is
standard: one is a split-key PRF and the other a PRF-ODH variant
(\cref{sec:mech:kem}). We now name each assumption precisely, and we added
the random-oracle proof from standard assumptions for that combiner's
X25519 leg.

\paragraph{Modelling of X25519.} The first write-up said the X-Wing models
make ``exactly the X-Wing paper's assumptions''. They do, up to one
difference: the models use a prime-order group and the paper a nominal
group. The review also noted that pyca/cryptography raises an error on an
all-zero X25519 output, while the models have no such failure; this led to
\cref{sec:lowcaught}.

\paragraph{Restrictions missing from the result statements.} The first
write-up disclosed two restrictions, in its modelling notes and in a list of
modelling choices that limit the statement. Keys of sessions with a
\emph{possible} honest partner were withheld from RevealKey, which the notes
called ``a small weakening of the adversary that the model does not
mechanically justify''. And ephemeral-key reveal was static. For an
initiator, a possible partner was a responder session with the same
transcript. For a responder, it was an initiator that had sent the same
$m_1$ and expected this responder, matched before the transcript was fixed.
The statements of the results, however, carried neither restriction; they
said that all four claims of the sketch were machine-checked. We removed the
first
restriction: every non-fresh session's key can now be revealed on request,
with partner-based freshness on the full transcript (\cref{sec:freshness}).
We also added the key-uniqueness query and guard mutants. All queries
remained provable. The claims now say that ephemeral reveal is static.

\paragraph{Evidence that was weaker than it looked.} The first version of the
$\ct$ mutant showed only that the full proof script aborted once $\ct$ was
removed from $\THH$. That shows the script depends on $\ct$, not that
agreement does. The mutant now runs an authentication-only proof, on the
real model as a control and on the mutant, and checks the agreement queries
directly (\cref{sec:mutations}). The review also asked for the composition
side conditions to be listed; they are in \cref{sec:status}.

\paragraph{The second review.} A second instance reviewed a draft of this
paper together with the repository. Besides many corrections to the draft,
it found four things in the repository. One was the missing factor in the
sketch's Theorem~1 (\cref{sec:sketchcaught}). Another was that the nine
mutants of the implementation fix existed only as a one-off run, and are now
a script in the artifact. A third was that the artifact commit was not on
the repository's main branch; the reproduction instructions now name a
branch that contains it. The fourth was that our claim that the
freshness mutants ``remove each of these guards in turn'' overstated their
coverage. The six freshness mutants then covered seven of the model's
twelve guards, two of them only together. We added mutants until every
guard had its own, writing the expected outcomes before running them. All
twenty-two tests behave as expected (\cref{sec:mutations}).

\subsection{In the implementation: low-order X25519 points}\label{sec:lowcaught}

X-Wing~\cite{XWingDraft} takes X25519 from RFC~7748~\cite[\S5]{RFC7748} and
applies no check to its output. Encapsulate and Decapsulate
(\S5.4--5.5 of the draft) pass the raw Diffie--Hellman value to the
combiner, and the draft gives Decapsulate no failure case. It does not
forbid a check either. RFC~7748 makes the all-zero check optional
(``MAY''~\cite[\S6.1]{RFC7748}), while RFC~9180 requires DHKEM recipients to
perform it (\S7.1.4)~\cite{RFC9180} and RFC~10024 requires it for hybrid key
agreement in TLS~1.3 (\S4.3)~\cite{RFC10024}. Our implementation computed
X25519 with pyca/cryptography (version 50.0.1), which raises
\texttt{ValueError} when the output is all-zero. The first review noted this
divergence from the models, and tests confirmed what it meant for X-Wing:
our decapsulation raised on any ciphertext whose X25519 part is a low-order
point, and our encapsulation raised on any such public key, where the
draft's algorithm derives a key.

Neither the models nor the known-answer tests could see this. The models
treat X25519 as a prime-order group, which has no low-order points. The
draft's test vectors use honestly generated values, which are never
low-order. The divergence did not affect secrecy. Honest parties never
produce these points, and whether a point is low-order is a function of
public data, so the error gave an adversary nothing it could not compute
itself. It was an interoperability difference: a peer that follows the
draft's algorithm derives a key where ours raised, so two implementations of
the same KEM disagreed about which inputs have a key.

The divergence was not new. The same behaviour had been reported for the
noble library's ML-KEM-768/X25519 KEM, whose X25519 implementation throws on
an all-zero output~\cite{Noble26}. Neither is the list of affected inputs,
which libraries such as Mbed~TLS check against~\cite{mbedtls}. What we add is
a fix that still raises every other error, a self-contained proof that the
list is complete, and tests that pin both.

\paragraph{The fix.} For clamped scalars, the X25519 output is all-zero
exactly when the peer's u-coordinate, decoded as RFC~7748 prescribes (top bit
masked, value reduced mod $p$), is one of five values. These are $0$, $1$,
$p-1$, and the u-coordinates of the two pairs of points of order~8.
\Cref{app:loworder} proves this: seven integers below $2^{255}$ represent
these residues, giving fourteen 32-byte encodings. The implementation now
catches pyca's error and returns the RFC~7748 value, 32 zero bytes, only if
the peer's input decodes to one of the five values; any other error is
raised again:
\begin{lstlisting}[language=Python,morekeywords={try,except,raise,def,return,if,in}]
def _x25519_dh(sk, peer_pk):
    priv = x25519.X25519PrivateKey.from_private_bytes(sk)
    peer = x25519.X25519PublicKey.from_public_bytes(peer_pk)
    try:
        return priv.exchange(peer)
    except ValueError:
        if _is_low_order(peer_pk):    # public input only
            return bytes(32)          # the RFC 7748 output
        raise
\end{lstlisting}
The ML-KEM-1024 profile keeps its check, following RFC~9180's rule for DHKEM
recipients~\cite[\S7.1.4]{RFC9180} in both directions. Its models do not
include that rejection either, but the same argument applies: it depends
only on public data.

\paragraph{Tests.} Sixty-three new tests cover the fix, and \cref{app:loworder}
describes them in detail. An independent RFC~7748 Montgomery ladder, checked
against RFC~7748's first test vector and the first iteration of its iterated
test and against pyca on random inputs, confirms that each of the fourteen
encodings gives an all-zero output. Counting roots over $\Fp$ confirms that
no other u-coordinate has order dividing~8. pyca rejects each encoding, so
the fallback is reachable, and it accepts random and near-miss inputs, so the
fallback is not taken spuriously. Decapsulation and encapsulation return the
draft's value, recomputed with \texttt{hashlib} rather than through the code
under test, and the ML-KEM-1024 profile still rejects. A script in the
artifact, \file{tests/low_order_mutants.py}, applies nine mutants of the fix
and checks that each makes some of these tests fail. The mutants remove the
fallback, catch every error, drop one point from the set, skip the top-bit
mask or the reduction mod $p$, revert either direction, return a wrong
constant, or apply the fallback to the RFC~9180 leg as well. The full test
suite (335 tests) passes.

\subsection{Lessons}\label{sec:lessons}

\paragraph{Mechanize the primitive lemmas, not only the protocol.} Two of
the three errors in the sketch were in ``secure if either component is''
lemmas, written as a minimum of two advantages without side assumptions.
Statements of that shape read as obviously true. In both cases the primary
source states the missing condition, the composite draft~\cite[\S1.3]{Composite}
and the X-Wing paper~\cite[Thm.~1]{XWing}, and the sketch's paraphrase
dropped it. Mechanization put it back as soon as the construction was
modelled byte-for-byte as specified: the composite's message representative,
and X-Wing's choice of combiner inputs.

\paragraph{Compare bounds term by term.} The mechanized Theorem~1 bound had
the factor $N_K$ that the sketch's lacked, and we did not notice until a
reviewer compared the two terms. A mechanized bound is evidence about the
pen-and-paper one only if someone reads it that way.

\paragraph{A checked proof does not check its description.} Most review
findings concerned what we said about the models: how tight the bound is,
what an assumption is called, which restrictions the claims carry, how much
the mutation tests cover. CryptoVerif checks none of these. Claims should be
read off the bound's terms and the library macros' definitions, and checked
against the models by someone who did not write them.

\paragraph{Make negative controls test the property, and count what they
cover.} A proof script that aborts after a mutation is weak evidence, since
scripts abort for many reasons. A mutant should check that the query it
targets is no longer proved, next to a control in which the same proof
succeeds on the unmutated model. Coverage should be checked against an
explicit list of what is meant to be covered, such as the model's guards.

\paragraph{Test at the abstraction boundary.} The implementation diverged
exactly where the models abstract (low-order points do not exist in a
prime-order group) and where the test vectors do not reach (they are honest
values). The review found the divergence by asking what the abstraction
leaves out, and tests at that boundary confirmed it.
